\chapter{Assisted Vaginal Delivery (Antenatal Use)}
\index{Assisted Vaginal Delivery (Antenatal Use)}

\section*{Definition and Overview}
Assisted vaginal delivery, also called instrumental delivery, is a birth in which the baby is helped out of the birth canal with the use of instruments, either a vacuum cup (ventouse) or forceps. It may be needed when the second stage of labour is prolonged, when the baby shows signs of distress, or when the mother is exhausted and unable to push effectively. This chapter is written for use during pregnancy, so that expectant mothers understand why assisted delivery happens, how it is performed, and what to expect, even though the decision to use instruments can only be made during labour itself.

\section*{Epidemiology}
Around 10 to 15 per cent of vaginal births involve an assisted delivery in many countries, although rates vary between hospitals and regions. Assisted delivery is more likely in first pregnancies, when epidural anaesthesia is used, when labour is induced, and when the baby is large or lying in an awkward position. It is an alternative to caesarean section in the second stage of labour, and the choice between them depends on the clinical situation, the skill of the operator, and the facilities available.

\section*{Aetiology and Pathophysiology}
Assisted vaginal delivery is considered when the natural second stage of labour is not progressing or when there is concern about the wellbeing of the baby or mother.
\begin{itemize}
    \item \textbf{Prolonged second stage:} the baby has not descended despite an adequate period of pushing
    \item \textbf{Fetal distress:} changes in the baby's heart-rate pattern suggest it may not tolerate further pushing
    \item \textbf{Maternal exhaustion:} the mother is too tired to push effectively
    \item \textbf{Maternal medical conditions:} where prolonged pushing is hazardous, such as severe heart disease, poorly controlled hypertension, or raised pressure inside the eye
    \item \textbf{Malposition:} the baby presenting face-up or in another position that makes spontaneous descent slower
    \item \textbf{Shortening the pushing phase:} in some circumstances, to reduce maternal effort and expedite birth
\end{itemize}

\section*{Clinical Features}
From the mother's perspective, an assisted delivery is recognised by the steps that accompany it:
\begin{itemize}
    \item A clear explanation of why assistance is needed and how it will be performed
    \item Placement of a vacuum cup on the baby's head, or forceps blades around the head
    \item Contractions and gentle traction timed with the mother's pushing
    \item An episiotomy, a small surgical cut, sometimes needed to make more room
    \item Increased monitoring of the baby's heart rate throughout the birth
    \item The presence of an experienced doctor or midwife who performs the procedure
\end{itemize}

\section*{Differential Diagnosis}
When the second stage of labour stalls, the possible reasons that must be distinguished include:
\begin{itemize}
    \item Poor or inefficient uterine contractions
    \item A large baby or a baby in a difficult position (malposition)
    \item A small or distorted maternal pelvis preventing descent
    \item Fetal distress requiring prompt delivery
    \item Maternal exhaustion or ineffective pushing technique
    \item The alternative of proceeding to caesarean section rather than assisted vaginal birth
\end{itemize}

\section*{When to Seek Medical Care}
During pregnancy and labour, seek help promptly if you have any concern about your baby's movements, if contractions become very intense, or if you are in labour and feel the urge to push before your midwife or doctor has assessed you.

\textbf{Contact your local emergency services (or attend an emergency department) for:}
\begin{itemize}
    \item Heavy vaginal bleeding in pregnancy or labour
    \item Severe abdominal pain, or a prolonged period without feeling the baby move
    \item A sudden gushing of fluid, or signs that the umbilical cord may have come down
    \item Seizures, a severe headache with visual disturbance, or collapse in pregnancy
\end{itemize}

\section*{Diagnosis and Investigations}
\begin{itemize}
    \item \textbf{Assessment in labour:} vaginal examination to confirm full cervical dilatation and to establish how low the baby's head has descended
    \item \textbf{Fetal monitoring:} cardiotocography (CTG) to assess the baby's heart-rate pattern and responses
    \item \textbf{Examination of the baby's position:} by abdominal palpation and internal examination, because position determines the choice of instrument
    \item \textbf{Maternal observations:} temperature, blood pressure, and pulse to detect fever, infection, or maternal distress
    \item \textbf{Discussion and consent:} assisted delivery requires the mother's understanding and agreement before it is performed
\end{itemize}

\section*{Management}
\begin{itemize}
    \item \textbf{Antenatal preparation:} birth-plan discussions, positioning, and pushing techniques taught in childbirth classes
    \item \textbf{Analgesia:} epidural or other pain relief is often already in place; additional local anaesthesia may be required for forceps
    \item \textbf{Ventouse:} a soft vacuum cup is applied to the baby's head and gentle traction is applied during contractions
    \item \textbf{Forceps:} shaped metal blades are placed around the baby's head to provide traction and correct rotation
    \item \textbf{Episiotomy:} a surgical cut to enlarge the vaginal opening when the perineum will not stretch sufficiently
    \item \textbf{After the birth:} mother and baby are observed closely, and any perineal tears or cuts are repaired
\end{itemize}

\section*{Complications}
\begin{itemize}
    \item \textbf{Maternal:} perineal tears, heavier bleeding after birth, bruising, and local discomfort
    \item \textbf{Future continence:} an increased chance of pelvic floor weakness and bladder or bowel symptoms in some women
    \item \textbf{Baby:} bruising or swelling of the scalp from the vacuum cup, temporary marks from forceps, or rarely a small bleed beneath the skull; serious injury is uncommon
    \item \textbf{Shoulder dystocia:} the baby's shoulders becoming stuck behind the pelvis, which is an obstetric emergency
    \item \textbf{Failed assisted delivery:} occasionally the baby must be delivered by caesarean section after all
\end{itemize}

\section*{Prognosis and Outcomes}
Most assisted vaginal deliveries result in a healthy mother and baby, and the procedure avoids caesarean section in many cases. Scalp bruising and perineal discomfort usually settle within days to weeks. Pelvic floor symptoms such as urinary leakage may persist in some women and can often be improved with pelvic floor physiotherapy. Most women who have had one assisted delivery are able to have a spontaneous vaginal birth in a subsequent pregnancy.

\section*{Prevention and Self-Care}
\begin{itemize}
    \item Attend antenatal education to learn about labour, positioning, and pushing
    \item Maintain good general health in pregnancy through regular antenatal care and a balanced diet
    \item Practise pelvic floor exercises during and after pregnancy
    \item Discuss your birth preferences with your midwife or obstetrician before the birth
    \item Ask about the reasons for, and alternatives to, any intervention during labour
    \item After the birth, attend follow-up and seek help early for pain, bleeding, or difficulty passing urine or opening the bowels
\end{itemize}

\section*{Sources and Further Reading}
The following organisations provide reliable, evidence-based information on assisted vaginal delivery for parents and healthcare students.

\begin{itemize}
    \item NHS. \textit{Assisted delivery (ventouse or forceps).} https://www.nhs.uk/pregnancy/labour-and-birth/what-happens/assisted-delivery/
    \item Royal College of Obstetricians and Gynaecologists (RCOG). \textit{Operative vaginal delivery patient information.} https://www.rcog.org.uk/
    \item MedlinePlus. \textit{Vaginal delivery with forceps or vacuum.} https://medlineplus.gov/ency/patientinstructions/000553.htm
    \item Mayo Clinic. \textit{Labour and delivery: assisted delivery.} https://www.mayoclinic.org/healthy-lifestyle/labor-and-delivery/
    \item World Health Organization. \textit{Maternal and newborn health.} https://www.who.int/health-topics/maternal-health
    \item Centers for Disease Control and Prevention (CDC). \textit{Assisted vaginal delivery information.} https://www.cdc.gov/reproductivehealth/maternalinfanthealth/
\end{itemize}
